\textbf{\textit{Soal. }}Given triangle $ABC$ with incenter $I$ and $A$-excenter $J$. Circle $\omega_b$ centered at point $O_b$ passes through point $B$ and is tangent to line $CI$ at point $I$. Circle $\omega_c$ with center $O_c$ passes through point $C$ and touches line $BI$ at point $I$. Let $O_bO_c$ and $IJ$ intersect at point $K$. Find the ratio $IK/KJ$.

\begin{proof}[\textbf{Solusi. }] (Friday, 28 October 2022) Misalkan $M_b$, $M_c$, dan $M$ berturut-turut adalah titik tengah $IB$, $IC$, dan $IJ$. Dari lemma terkenal, $M$ adalah titik pusat lingkaran $(BICJ)$ dan $M$ berada di lingkaran $(ABC)$. Oleh karena itu, $M_bM$ dan $M_cM$ adalah garis sumbu lingkaran $(BIC)$. Di lain sisi, kita juga punya $O_bM_b$ dan $O_cM_c$ berturut-turut adalah garis sumbu dari lingkaran $\omega_b$ dan $\omega_c$. Hal ini menyebabkan $M, O_b, M_b$ segaris dan $M, O_c, M_c$ segaris. Namun, sadari bahwa $MO_b \perp BI$ dan $O_cI \perp BI$ (karena $I$ merupakan titik singgung $\omega_c$) yang menyebabkan $MO_b \parallel O_cI$. Dengan cara serupa $MO_c \parallel O_bI$ juga. Oleh karena itu kita peroleh $MO_bIO_c$ adalah sebuah jajar genjang sehingga $K$ merupakan titik tengah dari $IM$ atau $IK = \dfrac{1}{2}IM = \dfrac{1}{4}IJ$. Ini berakibat $\dfrac{IK}{KJ} = \dfrac{\frac{1}{4}IJ}{\frac{3}{4}IJ}=\boxed{\dfrac{1}{3}}$.
% \begin{center}
    \begin{asy}
        import olympiad;
        import geometry;
        size(350);

        pair A, B, C;
        A = dir(125);
        B = dir(-150);
        C = dir(-30);
        dot("$A$", A, NW);
        dot("$B$", B, W);
        dot("$C$", C, E);
        draw(A--B--C--A);

        circle g = circumcircle(triangle(A,B,C));
        draw(g);

        pair I = incenter(triangle(A,B,C));
        dot("$I$", I, NE);
        
        pair intA[] = intersectionpoints(line(A,I), g);

        pair M = intA[0];
        pair J = rotate(180,M)*I;
        dot("$M$", M, SE);
        dot("$J$", J, SE);
        draw(circumcircle(triangle(B,I,C)), blue);

        draw(A--J);     
        draw(B--I, blue);
        draw(C--I, blue);

        line icperp = perpendicular(I, line(I,C));
        pair Mb = (I+B)/2;
        line Mbperp = perpendicular(Mb, line(I,B));
        pair Ob = intersectionpoint(icperp, Mbperp);
        dot("$M_b$", Mb, NW);
        dot("$O_b$", Ob, NE);
        pair Mbb = foot(Ob, I, M);
        pair Mbbb = rotate(180, Mbb)*I;
        draw(circumcircle(triangle(I,B,Mbbb)), red);
        
        line ibperp = perpendicular(I, line(I,B));
        pair Mc = (I+C)/2;
        line Mcperp = perpendicular(Mc, line(I,C));
        pair Oc = intersectionpoint(ibperp, Mcperp);
        dot("$M_c$", Mc, NE);
        dot("$O_c$", Oc, SE);
        pair Mcc = foot(Oc, I, M);
        pair Mccc = rotate(180, Mcc)*I;
        draw(circumcircle(triangle(I,C,Mccc)), red);

        draw(Mb--M, purple);
        draw(Mc--M, purple);
        draw(I--Ob--M--Oc--I, green);

        pair K = intersectionpoint(line(I,M),line(Ob,Oc));
        draw(Ob--Oc, green);
        dot("$K$", K, SW);

        draw(rightanglemark(Ob, I, C, 3), green);
        draw(rightanglemark(M, Mc, C, 2));
        draw(rightanglemark(Oc, I, B, 1.5), red);
        draw(rightanglemark(M, Mb, B, 2));
    \end{asy}
\end{center}
\begin{center}
    \begin{tikzpicture}[scale=3.2, dot/.style={circle, fill, inner sep=0.5pt, outer sep=0pt}]

        % Titik-titik utama segitiga ABC
        \tkzDefPoint(125:1){A}
        \tkzDefPoint(210:1){B}
        \tkzDefPoint(-30:1){C}

        % Lingkaran luar dan pusatnya
        \tkzDefTriangleCenter[circum](A,B,C)\tkzGetPoint{O}
        \tkzDrawCircle(O,A)

        % Incenter
        \tkzInCenter(A,B,C)\tkzGetPoint{I}

        % M titik tengah busur BC (perpotongan kedua AI dengan lingkaran luar)
        \tkzInterLC[common=A](A,I)(O,A)\tkzGetFirstPoint{M}

        % J refleksi I terhadap M
        \tkzDefPointBy[symmetry=center M](I)\tkzGetPoint{J}

        % Lingkaran (BIC)
        \tkzDefTriangleCenter[circum](B,I,C)\tkzGetPoint{Obic}
        \tkzDrawCircle[blue](Obic,B)

        \tkzDrawSegment(A,J)
        \tkzDrawSegments[blue](B,I C,I)

        % Konstruksi Ob: titik potong tegak lurus IC di I dan tegak lurus IB di titik tengah IB
        \tkzDefMidPoint(I,B)\tkzGetPoint{Mb}
        \tkzDefLine[orthogonal=through I](I,C)\tkzGetPoint{icp}
        \tkzDefLine[orthogonal=through Mb](I,B)\tkzGetPoint{mbp}
        \tkzInterLL(I,icp)(Mb,mbp)\tkzGetPoint{Ob}
        \tkzDefPointBy[projection=onto I--M](Ob)\tkzGetPoint{Mbb}
        \tkzDefPointBy[symmetry=center Mbb](I)\tkzGetPoint{Mbbb}
        \tkzDefTriangleCenter[circum](I,B,Mbbb)\tkzGetPoint{Oibm}
        \tkzDrawCircle[red](Oibm,I)

        % Konstruksi Oc: simetris dengan menukar B dan C
        \tkzDefMidPoint(I,C)\tkzGetPoint{Mc}
        \tkzDefLine[orthogonal=through I](I,B)\tkzGetPoint{ibp}
        \tkzDefLine[orthogonal=through Mc](I,C)\tkzGetPoint{mcp}
        \tkzInterLL(I,ibp)(Mc,mcp)\tkzGetPoint{Oc}
        \tkzDefPointBy[projection=onto I--M](Oc)\tkzGetPoint{Mcc}
        \tkzDefPointBy[symmetry=center Mcc](I)\tkzGetPoint{Mccc}
        \tkzDefTriangleCenter[circum](I,C,Mccc)\tkzGetPoint{Oicm}
        \tkzDrawCircle[red](Oicm,I)

        \tkzDrawSegments[purple](Mb,M Mc,M)
        \tkzDrawSegments[green!60!black](I,Ob Ob,M M,Oc Oc,I)
        \tkzDrawSegment[green!60!black](Ob,Oc)

        % K perpotongan IM dan ObOc
        \tkzInterLL(I,M)(Ob,Oc)\tkzGetPoint{K}

        % Tanda sudut siku-siku
        \tkzMarkRightAngle[green!60!black, size=0.0469](Ob,I,C)
        \tkzMarkRightAngle[red, size=0.0469](Oc,I,B)
        \tkzMarkRightAngle[size=0.0469](M,Mb,B)
        \tkzMarkRightAngle[size=0.0469](M,Mc,C)

        % Titik dan label
        \tkzDrawPoints[fill=black, size=3](A,B,C,I,M,J,Mb,Ob,Mc,Oc,K)
        \tkzLabelPoint[above left](A){$A$}
        \tkzLabelPoint[left](B){$B$}
        \tkzLabelPoint[right](C){$C$}
        \tkzLabelPoint[above right](I){$I$}
        \tkzLabelPoint[below right](M){$M$}
        \tkzLabelPoint[below](J){$J$}
        \tkzLabelPoint[above left](Mb){$M_b$}
        \tkzLabelPoint[above right](Ob){$O_b$}
        \tkzLabelPoint[above right](Mc){$M_c$}
        \tkzLabelPoint[below right](Oc){$O_c$}
        \tkzLabelPoint[below left](K){$K$}

    \end{tikzpicture}
\end{center}

\end{proof}